\scanpage{42}
% archive foliation: 41
D'après l'hypothèse \deleted{les $\varphi^{-1}(U_i)$ \uline{recouvrent} $T$,} pour tt $s\in S$, il existe un $i(s)$ tel que $\varphi(T_s)\subset U_{i(s)}$. Il est évident (préciser) que l'on peut trouver un voisinage $V_{(s)}$ de $s$ tel que $\varphi(T\vert V_{(s)})\subset U_i$. On en déduit un hom $V_{(s)}\longrightarrow H_i$, d'où $V_{(s)}\longrightarrow H$.
Je dois montrer que ça se recolle. C'est \unsure{strictement formel} ---.

\uline{Théorème} On suppose $T$ fini et loc. libre sur $S$,
\deleted{et $X/S$ quelconque. Pour que \unsure{$\operatorname{Hom}$}}
On suppose de plus que \deleted{\unsure{l'on ait dans}} \deleted{l'une des deux conditions suivantes\par (i) $T\longrightarrow S$ est \unsure{univ. injectif|purement injectif}\par (ii)}
Pour tout $\unsure{t\in S|s\in S}$, tt multisection de $X_s$ sur $T_s$, définie sur une extension finie de $k(s)$, soit contenue dans un ouvert affine de $X$.
\margin{$k(s)\qquad k(t)$}
\deleted{[\unsure{condition Hyp.}} \inserted{Alors} \deleted{\unsure{enfin si tt partie}}
$f_*(X/T)$ existe, et est réunion des $f_{i*}(U_i/T)$, où $U_i$ parcourent ts ouverts affines de $X$.
